\documentclass[12pt,a4paper]{article}
\usepackage[utf8]{inputenc}
\usepackage[spanish]{babel}
\usepackage{graphicx}
\usepackage{hyperref}
\usepackage{geometry}
\usepackage{fancyhdr}
\usepackage{listings}
\usepackage{xcolor}
\usepackage{titlesec}

% Configuración de página
\geometry{margin=2.5cm}

% Configuración de enlaces
\hypersetup{
    colorlinks=true,
    linkcolor=blue,
    filecolor=magenta,      
    urlcolor=cyan,
    pdftitle={Segunda Entrega - Todo App + Jenkins},
}

% Configuración de código
\lstset{
    basicstyle=\ttfamily\footnotesize,
    breaklines=true,
    frame=single,
    backgroundcolor=\color{gray!10},
    commentstyle=\color{green!50!black},
    keywordstyle=\color{blue},
    stringstyle=\color{red},
    showstringspaces=false
}

% Encabezado y pie de página
\pagestyle{fancy}
\fancyhf{}
\rhead{Segunda Entrega - Jenkins CI/CD}
\lhead{DevOps - Todo App}
\rfoot{Página \thepage}

\begin{document}

% ===============================
% PORTADA
% ===============================
\begin{titlepage}
    \centering
    \vspace{2cm}
    
    {\Huge\bfseries Segunda Entrega del Proyecto\par}
    \vspace{1cm}
    {\LARGE Todo App + Jenkins CI/CD\par}
    \vspace{0.5cm}
    {\large Integración Continua y Despliegue Automatizado\par}
    
    \vspace{3cm}
    
    {\Large\textbf{Integrantes:}\par}
    \vspace{0.5cm}
    {\large
        JAVIER ALBEIRO AGAMEZ SILGADO\\[0.3cm]
        OSCAR DUVAN ALVAREZ CARO\\[0.3cm]
        NEIDY MICHELLE CRUZ ROSALEZ\\[0.3cm]
        MICHELLE MOREA NINCO\\[0.3cm]
        LEANDRO NICOLAS VARGAS CUBIDES\\[0.3cm]
    }
    
    \vfill
    
    {\large\textbf{Fecha de entrega:} [FECHA AQUÍ]\par}
    \vspace{0.5cm}
    {\large\textbf{Repositorio GitHub:}\par}
    \vspace{0.3cm}
    {\small\url{https://github.com/LeandroNV/todo-app-fullstack}\par}
    
    \vspace{1cm}
\end{titlepage}

% ===============================
% ÍNDICE
% ===============================
\tableofcontents
\newpage

% ===============================
% INTRODUCCIÓN
% ===============================
\section{Introducción}

Esta segunda entrega del proyecto Todo App representa un avance significativo en la implementación de prácticas DevOps modernas, específicamente en el ámbito de la Integración Continua y el Despliegue Continuo (CI/CD).

\subsection{Conexión con la Primera Entrega}

En la primera entrega, se logró:
\begin{itemize}
    \item Desarrollar una aplicación full-stack funcional con Vue 3, Node.js y MongoDB
    \item Dockerizar completamente la aplicación con contenedores optimizados
    \item Implementar docker-compose para orquestación de servicios
    \item Establecer una arquitectura de microservicios robusta
\end{itemize}

\subsection{Objetivo de esta Segunda Entrega}

El objetivo principal de esta entrega es \textbf{automatizar el proceso de integración y despliegue} mediante Jenkins, transformando el flujo de trabajo manual en un pipeline automatizado que:

\begin{enumerate}
    \item \textbf{Detecte cambios automáticamente:} A través de webhooks de GitHub
    \item \textbf{Valide el código:} Ejecutando linters y análisis estático
    \item \textbf{Compile el proyecto:} Transpilando TypeScript a JavaScript
    \item \textbf{Construya imágenes Docker:} Generando contenedores optimizados
    \item \textbf{Ejecute pruebas:} Asegurando la calidad del código
    \item \textbf{Despliegue automáticamente:} Publicando la aplicación en producción
\end{enumerate}

\subsection{Beneficios Esperados}

La implementación de Jenkins en nuestro proyecto aporta:

\begin{itemize}
    \item \textbf{Automatización:} Eliminación de tareas manuales repetitivas
    \item \textbf{Consistencia:} Proceso de build y deploy estandarizado
    \item \textbf{Rapidez:} Despliegues frecuentes y confiables
    \item \textbf{Calidad:} Detección temprana de errores
    \item \textbf{Trazabilidad:} Historial completo de builds y despliegues
\end{itemize}

\newpage

% ===============================
% ARQUITECTURA DEL SISTEMA
% ===============================
\section{Arquitectura del Sistema}

\subsection{Componentes del Sistema}

El sistema integra cuatro componentes principales que trabajan en conjunto:

\begin{enumerate}
    \item \textbf{GitHub:} Sistema de control de versiones
    \item \textbf{Jenkins:} Servidor de integración continua
    \item \textbf{Docker:} Plataforma de contenedores
    \item \textbf{Aplicación:} Frontend + Backend + Base de datos
\end{enumerate}

\subsection{Flujo del Pipeline CI/CD}

El flujo completo del pipeline sigue esta secuencia:

\begin{lstlisting}[caption=Flujo del Pipeline CI/CD]
1. DESARROLLO
   └─> Desarrollador hace push a GitHub
   
2. TRIGGER
   └─> Webhook notifica a Jenkins
   
3. CHECKOUT
   └─> Jenkins descarga el código desde GitHub
   
4. INSTALACIÓN
   ├─> Instala dependencias del frontend
   └─> Instala dependencias del backend
   
5. ANÁLISIS
   ├─> Ejecuta linters (ESLint)
   └─> Verifica calidad de código
   
6. COMPILACIÓN
   ├─> Compila TypeScript frontend
   └─> Compila TypeScript backend
   
7. BUILD DOCKER
   ├─> Construye imagen frontend
   └─> Construye imagen backend
   
8. TESTS
   ├─> Tests unitarios
   └─> Tests de integración
   
9. PUSH
   └─> Sube imágenes a Docker Hub
   
10. DEPLOY
    └─> Despliega con docker-compose
    
11. VERIFICACIÓN
    └─> Healthcheck de servicios
\end{lstlisting}

\subsection{Integración de Componentes}

\textbf{GitHub → Jenkins:}
\begin{itemize}
    \item Webhook configurado para notificar en cada push
    \item Jenkins clona el repositorio automáticamente
    \item Credenciales gestionadas de forma segura
\end{itemize}

\textbf{Jenkins → Docker:}
\begin{itemize}
    \item Jenkins ejecuta comandos Docker directamente
    \item Construye imágenes desde Dockerfiles
    \item Gestiona contenedores con docker-compose
\end{itemize}

\textbf{Docker → Aplicación:}
\begin{itemize}
    \item Contenedores aislados para cada servicio
    \item Red personalizada para comunicación
    \item Volúmenes para persistencia de datos
\end{itemize}

\subsection{Repositorio del Proyecto}

\textbf{URL:} \url{https://github.com/LeandroNV/todo-app-fullstack}

\textbf{Estructura:}
\begin{lstlisting}
todo-app-fullstack/
├── Jenkinsfile              # Pipeline CI/CD
├── docker-compose.yml       # Orquestación
├── Dockerfile               # Imagen frontend
├── backend/
│   └── Dockerfile          # Imagen backend
└── src/                    # Código fuente
\end{lstlisting}

\newpage

% ===============================
% CONFIGURACIÓN TÉCNICA
% ===============================
\section{Configuración Técnica}

\subsection{Instalación de Jenkins}

\subsubsection{Requisitos Previos}

\begin{itemize}
    \item Java JDK 11 o superior
    \item Docker y Docker Compose instalados
    \item Acceso a GitHub con permisos de administrador
    \item 2GB RAM mínimo, 4GB recomendado
\end{itemize}

\subsubsection{Instalación en Docker}

\textbf{Paso 1: Crear directorio para Jenkins}
\begin{lstlisting}[language=bash]
mkdir -p jenkins_home
chmod 777 jenkins_home
\end{lstlisting}

\textbf{Paso 2: Ejecutar Jenkins en Docker}
\begin{lstlisting}[language=bash]
docker run -d \
  --name jenkins \
  -p 8080:8080 \
  -p 50000:50000 \
  -v jenkins_home:/var/jenkins_home \
  -v /var/run/docker.sock:/var/run/docker.sock \
  --restart=unless-stopped \
  jenkins/jenkins:lts
\end{lstlisting}

\textbf{Paso 3: Obtener contraseña inicial}
\begin{lstlisting}[language=bash]
docker exec jenkins cat /var/jenkins_home/secrets/initialAdminPassword
\end{lstlisting}

\textbf{Paso 4: Acceder a Jenkins}
\begin{itemize}
    \item Abrir navegador en \texttt{http://localhost:8080}
    \item Ingresar la contraseña inicial
    \item Instalar plugins sugeridos
\end{itemize}

\subsection{Configuración de Puertos}

\begin{table}[h]
\centering
\begin{tabular}{|l|l|l|}
\hline
\textbf{Servicio} & \textbf{Puerto} & \textbf{Descripción} \\
\hline
Jenkins Web UI & 8080 & Interfaz de administración \\
Jenkins Agent & 50000 & Comunicación con agentes \\
Frontend & 80 & Aplicación Vue \\
Backend API & 3000 & API REST \\
MongoDB & 27017 & Base de datos \\
\hline
\end{tabular}
\caption{Configuración de puertos del sistema}
\end{table}

\subsection{Plugins Instalados}

\begin{enumerate}
    \item \textbf{Git Plugin:} Integración con GitHub
    \item \textbf{Docker Plugin:} Gestión de contenedores
    \item \textbf{Docker Pipeline:} Comandos Docker en pipeline
    \item \textbf{GitHub Integration:} Webhooks y notificaciones
    \item \textbf{NodeJS Plugin:} Soporte para Node.js
    \item \textbf{Timestamper:} Marcas de tiempo en logs
    \item \textbf{Workspace Cleanup:} Limpieza automática
\end{enumerate}

\textbf{Instalación de plugins:}
\begin{enumerate}
    \item Ir a \texttt{Manage Jenkins → Manage Plugins}
    \item Buscar cada plugin en la pestaña "Available"
    \item Marcar e instalar sin reiniciar
    \item Reiniciar Jenkins cuando sea necesario
\end{enumerate}

\subsection{Conexión con GitHub}

\subsubsection{Configurar Credenciales}

\textbf{Paso 1: Generar Personal Access Token en GitHub}
\begin{enumerate}
    \item Ir a GitHub → Settings → Developer settings
    \item Personal access tokens → Generate new token
    \item Permisos: \texttt{repo}, \texttt{admin:repo\_hook}
    \item Copiar el token generado
\end{enumerate}

\textbf{Paso 2: Agregar credenciales en Jenkins}
\begin{enumerate}
    \item Jenkins → Manage Jenkins → Manage Credentials
    \item Global credentials → Add Credentials
    \item Kind: "Username with password"
    \item Username: tu usuario de GitHub
    \item Password: el Personal Access Token
    \item ID: \texttt{github-credentials}
\end{enumerate}

\subsubsection{Configurar Webhook}

\textbf{En GitHub:}
\begin{enumerate}
    \item Ir al repositorio → Settings → Webhooks
    \item Add webhook
    \item Payload URL: \texttt{http://TU\_JENKINS:8080/github-webhook/}
    \item Content type: \texttt{application/json}
    \item Events: "Just the push event"
    \item Active: marcado
\end{enumerate}

\subsection{Configuración de Docker en Jenkins}

\textbf{Dar permisos a Jenkins para usar Docker:}
\begin{lstlisting}[language=bash]
# Entrar al contenedor de Jenkins
docker exec -u root -it jenkins bash

# Instalar Docker CLI
apt-get update
apt-get install -y docker.io

# Agregar usuario jenkins al grupo docker
usermod -aG docker jenkins

# Salir y reiniciar Jenkins
exit
docker restart jenkins
\end{lstlisting}

\textbf{Verificar instalación:}
\begin{lstlisting}[language=bash]
docker exec jenkins docker --version
docker exec jenkins docker ps
\end{lstlisting}

\newpage

% ===============================
% JENKINSFILE
% ===============================
\section{Jenkinsfile}

El Jenkinsfile define el pipeline completo de CI/CD. A continuación se describe cada etapa:

\subsection{Estructura General}

\begin{lstlisting}[language=Java, caption=Estructura del Jenkinsfile]
pipeline {
    agent any
    environment { /* Variables */ }
    options { /* Configuraciones */ }
    stages { /* Etapas */ }
    post { /* Acciones finales */ }
}
\end{lstlisting}

\subsection{Variables de Entorno}

\begin{lstlisting}[language=Java]
environment {
    DOCKER_REGISTRY = 'docker.io'
    DOCKER_CREDENTIALS_ID = 'dockerhub-credentials'
    IMAGE_TAG = "${BUILD_NUMBER}"
    FRONTEND_IMAGE = "leandronv/todo-frontend"
    BACKEND_IMAGE = "leandronv/todo-backend"
}
\end{lstlisting}

\textbf{Descripción:}
\begin{itemize}
    \item \texttt{IMAGE\_TAG}: Versiona las imágenes con el número de build
    \item \texttt{FRONTEND\_IMAGE/BACKEND\_IMAGE}: Nombres de las imágenes en Docker Hub
    \item \texttt{DOCKER\_CREDENTIALS\_ID}: ID de credenciales para Docker Hub
\end{itemize}

\subsection{Etapa 1: Checkout}

\begin{lstlisting}[language=Java]
stage('Checkout') {
    steps {
        echo 'Descargando codigo desde GitHub...'
        checkout scm
    }
}
\end{lstlisting}

\textbf{Objetivo:} Descargar el código fuente desde GitHub.

\textbf{Resultado:} Código disponible en el workspace de Jenkins.

\subsection{Etapa 2: Verificar Dependencias}

Ejecuta verificaciones en paralelo para optimizar tiempo:

\begin{lstlisting}[language=Java]
parallel {
    stage('Verificar Node.js') {
        steps { sh 'node --version' }
    }
    stage('Verificar Docker') {
        steps { sh 'docker --version' }
    }
}
\end{lstlisting}

\textbf{Objetivo:} Asegurar que todas las herramientas necesarias estén instaladas.

\subsection{Etapa 3: Instalar Dependencias}

Instala dependencias de frontend y backend en paralelo:

\begin{lstlisting}[language=Java]
parallel {
    stage('Frontend Dependencies') {
        steps { sh 'npm install --legacy-peer-deps' }
    }
    stage('Backend Dependencies') {
        steps {
            dir('backend') { sh 'npm install' }
        }
    }
}
\end{lstlisting}

\textbf{Herramientas:} npm (Node Package Manager)

\textbf{Resultado:} Carpetas \texttt{node\_modules} creadas con todas las dependencias.

\subsection{Etapa 4: Análisis de Código}

Ejecuta linters para verificar calidad del código:

\begin{lstlisting}[language=Java]
stage('Lint Frontend') {
    steps {
        catchError(buildResult: 'SUCCESS', stageResult: 'UNSTABLE') {
            sh 'npm run lint || true'
        }
    }
}
\end{lstlisting}

\textbf{Herramientas:} ESLint, Prettier

\textbf{Resultado:} Reportes de problemas de estilo y errores potenciales.

\subsection{Etapa 5: Compilar TypeScript}

Transpila TypeScript a JavaScript:

\begin{lstlisting}[language=Java]
stage('Build Frontend') {
    steps { sh 'npm run build' }
}
stage('Build Backend') {
    steps {
        dir('backend') { sh 'npm run build' }
    }
}
\end{lstlisting}

\textbf{Herramientas:} TypeScript Compiler (tsc), Vite

\textbf{Resultado:} Carpeta \texttt{dist/} con código compilado.

\subsection{Etapa 6: Construir Imágenes Docker}

Crea imágenes Docker optimizadas:

\begin{lstlisting}[language=Java]
stage('Build Frontend Image') {
    steps {
        sh """
            docker build -t ${FRONTEND_IMAGE}:${IMAGE_TAG} \
                         -t ${FRONTEND_IMAGE}:latest .
        """
    }
}
\end{lstlisting}

\textbf{Herramientas:} Docker Engine

\textbf{Resultado:} Imágenes Docker versionadas y etiquetadas como \texttt{latest}.

\subsection{Etapa 7: Tests}

Ejecuta suite de pruebas:

\begin{lstlisting}[language=Java]
stage('Test Backend') {
    steps {
        dir('backend') {
            catchError(buildResult: 'SUCCESS') {
                sh 'npm test || echo "Tests no implementados"'
            }
        }
    }
}
\end{lstlisting}

\textbf{Herramientas:} Jest, Vitest

\textbf{Resultado:} Reportes de cobertura y tests exitosos/fallidos.

\subsection{Etapa 8: Push a Docker Hub}

Sube imágenes al registro:

\begin{lstlisting}[language=Java]
stage('Push a Docker Hub') {
    when { branch 'main' }
    steps {
        script {
            docker.withRegistry('https://index.docker.io/v1/', 
                              DOCKER_CREDENTIALS_ID) {
                sh """
                    docker push ${FRONTEND_IMAGE}:${IMAGE_TAG}
                    docker push ${FRONTEND_IMAGE}:latest
                """
            }
        }
    }
}
\end{lstlisting}

\textbf{Condición:} Solo se ejecuta en rama \texttt{main}.

\textbf{Resultado:} Imágenes disponibles en Docker Hub para despliegue.

\subsection{Etapa 9: Desplegar}

Despliega la aplicación con docker-compose:

\begin{lstlisting}[language=Java]
stage('Desplegar con Docker Compose') {
    when { branch 'main' }
    steps {
        sh '''
            docker compose down || true
            docker compose up -d
            sleep 30
            docker compose ps
        '''
    }
}
\end{lstlisting}

\textbf{Resultado:} Aplicación desplegada y funcionando.

\subsection{Etapa 10: Verificación}

Verifica que los servicios estén funcionando:

\begin{lstlisting}[language=Java]
stage('Verificacion de Salud') {
    steps {
        retry(3) {
            sh '''
                curl -f http://localhost:3000/health
                curl -f http://localhost/
            '''
        }
    }
}
\end{lstlisting}

\textbf{Resultado:} Confirmación de que todos los servicios están operativos.

\subsection{Post-Actions}

Acciones que se ejecutan al final:

\begin{lstlisting}[language=Java]
post {
    success {
        echo 'Pipeline ejecutado exitosamente!'
    }
    failure {
        echo 'Pipeline fallo!'
    }
    always {
        sh 'docker image prune -f || true'
    }
}
\end{lstlisting}

\newpage

% ===============================
% EVIDENCIAS
% ===============================
\section{Evidencias}

Esta sección debe incluir capturas de pantalla que demuestren el funcionamiento correcto del pipeline y Jenkins.

\subsection{Dashboard de Jenkins}

\textbf{Descripción esperada:} Captura del dashboard principal de Jenkins mostrando:
\begin{itemize}
    \item Lista de proyectos/jobs configurados
    \item Estado de los últimos builds (verde=éxito, rojo=fallo)
    \item Gráfico de tendencia de builds
    \item Panel de actividad reciente
\end{itemize}

\textbf{Elementos a mostrar:}
\begin{itemize}
    \item Job "todo-app-fullstack" visible
    \item Historial de builds (#1, #2, #3, etc.)
    \item Indicadores de estado (Weather icons)
\end{itemize}

\subsection{Configuración del Pipeline}

\textbf{Descripción esperada:} Captura de la configuración del job mostrando:
\begin{itemize}
    \item Tipo: Pipeline
    \item Definición: Pipeline script from SCM
    \item SCM: Git
    \item Repository URL: \url{https://github.com/LeandroNV/todo-app-fullstack}
    \item Script Path: Jenkinsfile
\end{itemize}

\subsection{Ejecución Completa del Pipeline}

\textbf{Descripción esperada:} Captura de la vista de stages mostrando:
\begin{itemize}
    \item Todas las etapas del pipeline visibles
    \item Tiempo de ejecución de cada stage
    \item Estados: Success (verde), Unstable (amarillo)
    \item Duración total del pipeline
\end{itemize}

\textbf{Stages a mostrar:}
\begin{enumerate}
    \item Checkout - XX seg
    \item Verificar Dependencias - XX seg
    \item Instalar Dependencias - XX seg
    \item Análisis de Código - XX seg
    \item Compilar TypeScript - XX seg
    \item Construir Imágenes Docker - XX seg
    \item Tests - XX seg
    \item Push a Docker Hub - XX seg
    \item Desplegar - XX seg
    \item Verificación de Salud - XX seg
\end{enumerate}

\subsection{Logs de Consola}

\textbf{Descripción esperada:} Captura de los logs mostrando:
\begin{itemize}
    \item Inicio del build con información del commit
    \item Salida de cada comando ejecutado
    \item Mensajes de éxito/error
    \item Resumen final del build
\end{itemize}

\textbf{Fragmentos importantes a capturar:}
\begin{lstlisting}
[Pipeline] echo
Descargando codigo desde GitHub...
[Pipeline] checkout
Cloning repository...
Commit: feat: Add Jenkinsfile for CI/CD
Author: Leandro Vargas

[Pipeline] echo
Instalando dependencias del frontend...
added 847 packages in 23.4s

[Pipeline] echo
Construyendo imagen Docker del frontend...
Step 1/10 : FROM node:20-alpine AS build
Successfully built abc123def456
Successfully tagged leandronv/todo-frontend:latest

[Pipeline] echo
Pipeline ejecutado exitosamente!
Duracion: 3m 45s
\end{lstlisting}

\subsection{Contenedores Desplegados}

\textbf{Descripción esperada:} Captura mostrando:
\begin{lstlisting}[language=bash]
$ docker compose ps

NAME              STATUS         PORTS
todo-frontend     Up 5 minutes   0.0.0.0:80->80/tcp (healthy)
todo-backend      Up 5 minutes   0.0.0.0:3000->3000/tcp (healthy)
todo-mongodb      Up 5 minutes   0.0.0.0:27017->27017/tcp (healthy)
\end{lstlisting}

\subsection{Aplicación Funcionando}

\textbf{Descripción esperada:} Captura del navegador mostrando:
\begin{itemize}
    \item URL: \texttt{http://localhost}
    \item Aplicación completamente funcional
    \item Indicador "Conectado" en verde
    \item Dashboard con estadísticas
    \item Formulario de tareas visible
\end{itemize}

\subsection{Healthcheck}

\textbf{Descripción esperada:} Captura o output de:
\begin{lstlisting}[language=bash]
$ curl http://localhost:3000/health

{
  "status": "ok",
  "timestamp": "2025-11-XX...",
  "uptime": 305.234
}
\end{lstlisting}

\subsection{Docker Hub}

\textbf{Descripción esperada:} Captura de Docker Hub mostrando:
\begin{itemize}
    \item Repositorios: \texttt{leandronv/todo-frontend} y \texttt{leandronv/todo-backend}
    \item Tags: \texttt{latest}, \texttt{1}, \texttt{2}, etc.
    \item Fecha de última actualización
    \item Tamaño de las imágenes
\end{itemize}

\newpage

% ===============================
% PROBLEMAS Y SOLUCIONES
% ===============================
\section{Problemas y Soluciones}

Durante la implementación de Jenkins, se encontraron varios desafíos técnicos:

\subsection{Permisos de Docker}

\textbf{Problema:} Jenkins no podía ejecutar comandos Docker debido a permisos insuficientes.

\textbf{Error recibido:}
\begin{lstlisting}
Got permission denied while trying to connect to the Docker daemon socket
\end{lstlisting}

\textbf{Solución:}
\begin{enumerate}
    \item Montar el socket de Docker en el contenedor de Jenkins:
    \begin{lstlisting}[language=bash]
-v /var/run/docker.sock:/var/run/docker.sock
    \end{lstlisting}
    
    \item Instalar Docker CLI dentro del contenedor de Jenkins
    
    \item Agregar el usuario jenkins al grupo docker:
    \begin{lstlisting}[language=bash]
docker exec -u root jenkins usermod -aG docker jenkins
docker restart jenkins
    \end{lstlisting}
\end{enumerate}

\subsection{Conflicto de Puertos}

\textbf{Problema:} El puerto 8080 ya estaba en uso por otro servicio.

\textbf{Solución:}
\begin{itemize}
    \item Identificar el proceso usando el puerto:
    \begin{lstlisting}[language=bash]
netstat -ano | findstr :8080
    \end{lstlisting}
    
    \item Cambiar el puerto de Jenkins a 8081:
    \begin{lstlisting}[language=bash]
docker run -d -p 8081:8080 -p 50000:50000 ...
    \end{lstlisting}
\end{itemize}

\subsection{Webhook de GitHub no funcionaba}

\textbf{Problema:} Jenkins no recibía notificaciones de GitHub en cambios al repositorio.

\textbf{Causa:} Jenkins corriendo en localhost no es accesible desde GitHub.

\textbf{Soluciones implementadas:}
\begin{enumerate}
    \item \textbf{Opción 1 - Ngrok:} Crear túnel público temporal
    \begin{lstlisting}[language=bash]
ngrok http 8080
# URL: https://abc123.ngrok.io
    \end{lstlisting}
    
    \item \textbf{Opción 2 - Polling:} Configurar Jenkins para verificar cambios periódicamente
    \begin{lstlisting}
Poll SCM: H/5 * * * *  # Cada 5 minutos
    \end{lstlisting}
\end{enumerate}

\subsection{Dependencias de npm con errores}

\textbf{Problema:} Conflictos de dependencias al ejecutar \texttt{npm install}.

\textbf{Error:}
\begin{lstlisting}
ERESOLVE unable to resolve dependency tree
\end{lstlisting}

\textbf{Solución:}
\begin{lstlisting}[language=bash]
npm install --legacy-peer-deps
\end{lstlisting}

\subsection{Build de Docker muy lento}

\textbf{Problema:} La construcción de imágenes Docker tomaba más de 10 minutos.

\textbf{Soluciones aplicadas:}
\begin{enumerate}
    \item Implementar multi-stage builds para reducir tamaño
    \item Optimizar el orden de capas en Dockerfile
    \item Usar \texttt{.dockerignore} para excluir archivos innecesarios
    \item Cachear capas de dependencias
\end{enumerate}

\textbf{Resultado:} Tiempo de build reducido a 3-4 minutos.

\subsection{Credenciales de Docker Hub}

\textbf{Problema:} Jenkins no podía autenticarse con Docker Hub para push.

\textbf{Solución:}
\begin{enumerate}
    \item Crear credenciales en Jenkins:
    \begin{itemize}
        \item Manage Jenkins → Manage Credentials
        \item Kind: Username with password
        \item Username: usuario de Docker Hub
        \item Password: Personal Access Token de Docker Hub
        \item ID: \texttt{dockerhub-credentials}
    \end{itemize}
    
    \item Usar en Jenkinsfile:
    \begin{lstlisting}[language=Java]
docker.withRegistry('https://index.docker.io/v1/', 
                   'dockerhub-credentials') {
    sh "docker push ${IMAGE}"
}
    \end{lstlisting}
\end{enumerate}

\subsection{Pipeline fallaba en stage de Tests}

\textbf{Problema:} Los tests aún no están implementados, causando que el pipeline falle.

\textbf{Solución:} Usar \texttt{catchError} para que fallas en tests no detengan el pipeline:

\begin{lstlisting}[language=Java]
catchError(buildResult: 'SUCCESS', stageResult: 'UNSTABLE') {
    sh 'npm test || echo "Tests no implementados aun"'
}
\end{lstlisting}

\newpage

% ===============================
% PRÓXIMOS PASOS
% ===============================
\section{Próximos Pasos}

Para la siguiente fase del proyecto, se planean las siguientes mejoras:

\subsection{Pruebas Automatizadas}

\subsubsection{Tests Unitarios}

\begin{itemize}
    \item \textbf{Frontend:} Implementar tests con Vitest
    \begin{itemize}
        \item Tests de componentes Vue
        \item Tests de stores (Pinia)
        \item Tests de servicios API
        \item Objetivo: 80\% de cobertura
    \end{itemize}
    
    \item \textbf{Backend:} Implementar tests con Jest
    \begin{itemize}
        \item Tests de controladores
        \item Tests de modelos Mongoose
        \item Tests de middleware
        \item Mocks de base de datos
    \end{itemize}
\end{itemize}

\subsubsection{Tests de Integración}

\begin{itemize}
    \item Usar Supertest para endpoints API
    \item Validar flujos completos CRUD
    \item Probar autenticación y autorización
    \item Tests con base de datos de prueba
\end{itemize}

\subsubsection{Tests End-to-End}

\begin{itemize}
    \item Implementar Playwright o Cypress
    \item Simular interacciones de usuario
    \item Tests de flujos críticos
    \item Capturas de pantalla en fallos
\end{itemize}

\subsection{Análisis de Cobertura de Código}

\begin{itemize}
    \item Integrar Codecov o SonarQube
    \item Reportes de cobertura en cada build
    \item Umbrales mínimos de cobertura
    \item Visualización de código no testeado
\end{itemize}

\subsection{Análisis de Seguridad}

\begin{itemize}
    \item \textbf{Trivy:} Escaneo de vulnerabilidades en imágenes Docker
    \item \textbf{npm audit:} Análisis de dependencias vulnerables
    \item \textbf{SAST:} Análisis estático de código
    \item \textbf{Secret scanning:} Detección de credenciales en código
\end{itemize}

\subsection{Notificaciones}

\begin{itemize}
    \item Integración con Slack para notificaciones de builds
    \item Emails en fallos críticos
    \item Webhooks a sistemas de monitoreo
    \item Notificaciones por commit
\end{itemize}

\subsection{Ambientes Múltiples}

\begin{itemize}
    \item \textbf{Development:} Rama develop
    \item \textbf{Staging:} Rama staging con datos de prueba
    \item \textbf{Production:} Rama main con despliegue controlado
    \item Pipelines específicos por ambiente
\end{itemize}

\subsection{Blue-Green Deployment}

\begin{itemize}
    \item Mantener dos ambientes idénticos
    \item Despliegue sin downtime
    \item Rollback instantáneo en caso de problemas
    \item Pruebas en ambiente blue antes de switch
\end{itemize}

\subsection{Monitoreo y Observabilidad}

\begin{itemize}
    \item \textbf{Prometheus:} Recolección de métricas
    \item \textbf{Grafana:} Dashboards de visualización
    \item \textbf{ELK Stack:} Centralización de logs
    \item Alertas configuradas para errores
\end{itemize}

\subsection{Performance Testing}

\begin{itemize}
    \item Tests de carga con K6 o Artillery
    \item Métricas de performance en cada build
    \item Detección de regresiones de performance
    \item Reportes de bottlenecks
\end{itemize}

\newpage

% ===============================
% CONCLUSIONES
% ===============================
\section{Conclusiones}

\subsection{Logros Alcanzados}

En esta segunda entrega se lograron objetivos significativos:

\begin{enumerate}
    \item \textbf{Automatización completa del pipeline:} Se eliminó la necesidad de builds y despliegues manuales, reduciendo errores humanos y tiempo de despliegue.
    
    \item \textbf{Integración exitosa con GitHub:} El sistema detecta automáticamente cambios en el repositorio y ejecuta el pipeline completo sin intervención manual.
    
    \item \textbf{Dockerización optimizada en CI:} Las imágenes Docker se construyen, versionan y publican automáticamente, manteniendo un registro histórico completo.
    
    \item \textbf{Pipeline robusto y confiable:} Con manejo de errores, reintentos y verificaciones de salud, el pipeline es resiliente y proporciona feedback claro.
    
    \item \textbf{Visibilidad completa:} Jenkins proporciona historial detallado de builds, logs y métricas, facilitando el debugging y análisis.
\end{enumerate}

\subsection{Aprendizajes Clave}

\subsubsection{Aspectos Técnicos}

\begin{itemize}
    \item \textbf{Configuración de Jenkins:} Aprendimos a instalar, configurar y gestionar Jenkins en Docker, incluyendo plugins y credenciales.
    
    \item \textbf{Sintaxis de Jenkinsfile:} Dominio del DSL de Jenkins para crear pipelines declarativos con stages paralelos y post-actions.
    
    \item \textbf{Integración Docker-Jenkins:} Comprendimos cómo Jenkins puede construir y gestionar contenedores Docker desde el pipeline.
    
    \item \textbf{Webhooks y automatización:} Implementamos comunicación bidireccional entre GitHub y Jenkins para triggers automáticos.
\end{itemize}

\subsubsection{Prácticas DevOps}

\begin{itemize}
    \item \textbf{Continuous Integration:} La importancia de validar código frecuentemente para detectar errores temprano.
    
    \item \textbf{Continuous Deployment:} Beneficios de despliegues automáticos y frecuentes para reducir riesgos.
    
    \item \textbf{Infrastructure as Code:} Jenkinsfile y docker-compose como código versionado facilitan reproducibilidad.
    
    \item \textbf{Feedback rápido:} Pipelines optimizados proporcionan feedback en minutos, acelerando el desarrollo.
\end{itemize}

\subsection{Relevancia de Jenkins}

Jenkins se ha consolidado como una herramienta fundamental en proyectos de desarrollo moderno por:

\begin{enumerate}
    \item \textbf{Código de calidad:} Al ejecutar linters y tests automáticamente, asegura que solo código validado llegue a producción.
    
    \item \textbf{Reducción de costos:} Automatización elimina tiempo dedicado a tareas repetitivas, permitiendo enfoque en desarrollo.
    
    \item \textbf{Despliegues confiables:} Proceso estandarizado reduce errores y permite rollbacks rápidos.
    
    \item \textbf{Colaboración mejorada:} Todo el equipo puede ver el estado del proyecto y acceder a builds históricos.
    
    \item \textbf{Escalabilidad:} A medida que el proyecto crece, Jenkins puede manejar pipelines más complejos con múltiples ambientes.
\end{enumerate}

\subsection{Impacto en el Proyecto}

La implementación de Jenkins ha transformado nuestro flujo de trabajo:

\begin{itemize}
    \item \textbf{Antes:} Builds manuales, despliegues inconsistentes, errores frecuentes, proceso lento.
    
    \item \textbf{Ahora:} Automatización completa, despliegues en minutos, proceso repetible, calidad asegurada.
\end{itemize}

\subsection{Reflexión Final}

Este proyecto ha demostrado que Jenkins no es solo una herramienta de automatización, sino un componente esencial de una cultura DevOps madura. La capacidad de integrar, validar y desplegar código de forma continua y confiable es fundamental en el desarrollo de software moderno.

La experiencia adquirida en esta entrega nos prepara para implementar prácticas más avanzadas como tests automatizados, análisis de seguridad y monitoreo continuo, consolidando un pipeline de CI/CD de nivel empresarial.

\newpage

% ===============================
% REFERENCIAS
% ===============================
\section{Referencias}

\begin{enumerate}
    \item Jenkins Official Documentation. \textit{Jenkins User Documentation}. \\
    \url{https://www.jenkins.io/doc/}
    
    \item Docker Documentation. \textit{Docker Documentation}. \\
    \url{https://docs.docker.com/}
    
    \item GitHub Docs. \textit{GitHub Webhooks Documentation}. \\
    \url{https://docs.github.com/en/webhooks}
    
    \item CloudBees. \textit{Jenkins Pipeline Syntax}. \\
    \url{https://www.jenkins.io/doc/book/pipeline/syntax/}
    
    \item Docker Hub. \textit{Jenkins Official Image}. \\
    \url{https://hub.docker.com/r/jenkins/jenkins}
    
    \item Atlassian. \textit{CI/CD Best Practices}. \\
    \url{https://www.atlassian.com/continuous-delivery/principles/continuous-integration-vs-delivery-vs-deployment}
    
    \item Mozilla Developer Network. \textit{HTTP Status Codes}. \\
    \url{https://developer.mozilla.org/en-US/docs/Web/HTTP/Status}
    
    \item Node.js Documentation. \textit{Node.js Official Docs}. \\
    \url{https://nodejs.org/docs/}
    
    \item Vue.js Documentation. \textit{Vue.js Guide}. \\
    \url{https://vuejs.org/guide/}
    
    \item MongoDB Documentation. \textit{MongoDB Manual}. \\
    \url{https://www.mongodb.com/docs/}
    
    \item Fowler, Martin. \textit{Continuous Integration}. \\
    \url{https://martinfowler.com/articles/continuousIntegration.html}
    
    \item GitHub Learning Lab. \textit{Introduction to GitHub}. \\
    \url{https://lab.github.com/}
    
    \item YouTube - TechWorld with Nana. \textit{Jenkins Tutorial for Beginners}. \\
    \url{https://www.youtube.com/watch?v=7KCS70sCoK0}
    
    \item freeCodeCamp. \textit{Docker Tutorial for Beginners}. \\
    \url{https://www.freecodecamp.org/news/the-docker-handbook/}
    
    \item Stack Overflow. \textit{Jenkins and Docker Integration Questions}. \\
    \url{https://stackoverflow.com/questions/tagged/jenkins+docker}
\end{enumerate}

\end{document}

